%! AP Statistics — Unit 1, Lesson 1.3 — Warm-Up ANSWER KEY
\documentclass[10pt]{article}
\usepackage{apstats-article}
\usepackage{apstats-key}

%=============================================================================
\begin{document}

\pageheader{Unit 1, Lesson 1.3}{Warm-Up --- Answer Key}
\begin{tabularx}{\linewidth}{X X X}
Name:\;\ans{Key} & Date:\; & Period:\;
\end{tabularx}

\vspace{0.18in}

\noindent Below is a list of 20 students' favorite school subjects.

\vspace{0.1in}

\begin{center}
\small
\begin{tabular}{*{10}{c}}
Math & Science & English & Math & Art \\
Science & Math & History & English & Math \\
Art & Science & Math & English & Science \\
History & Math & Art & Science & English \\
\end{tabular}
\end{center}

\vspace{0.18in}

\begin{enumerate}

  \item Can you tell which subject is most popular by just scanning the list above?
        Why or why not?\\[4pt]
        \ans{It is difficult to tell quickly because the data are unsorted. Tallying organizes
        the counts and makes the most popular category apparent.}

  \item Frequency table:

\vspace{0.1in}
\renewcommand{\arraystretch}{1.8}
\begin{tabularx}{\linewidth}{@{} >{\bfseries}p{0.22\linewidth}
                                    C{0.26\linewidth} C{0.2\linewidth} @{}}
\rowcolor{sky}
\textbf{Subject} & \textbf{Tally} & \textbf{Frequency} \\
\midrule
Art     & \ans{|||} & \ans{3} \\
English & \ans{||||} & \ans{4} \\
History & \ans{||} & \ans{2} \\
Math    & \ans{||||||} & \ans{6} \\
Science & \ans{|||||} & \ans{5} \\
\midrule
\rowcolor{sky}\textbf{Total} & & \ans{20} \\
\end{tabularx}

  \item Does the total match? \ans{Yes, 20.}\quad
        Why check? \ans{If the total doesn't equal $n$, a category was missed or
        double-counted --- the table is wrong.}

  \item Variable type: \ans{Categorical}\quad
        Meaningful to average? \ans{No. ``Favorite subject'' is a label, not a
        number. Averaging labels produces a meaningless result.}

\end{enumerate}

\vspace{0.14in}

\begin{tcolorbox}[colback=greenbg, colframe=greenacc, boxrule=0.8pt, arc=2mm,
    left=4mm, right=4mm, top=3mm, bottom=3mm]
\small
\begin{itemize}[leftmargin=1.3em, itemsep=6pt]
  \item The table you just built is a \textbf{frequency table}.
  \item Today we extend this to a \textbf{relative frequency table} by dividing each
        count by the total.
  \item Describing how often each category occurs = describing the
        \textbf{distribution} of a categorical variable.
\end{itemize}
\end{tcolorbox}

\end{document}
